% =============================================================================
% Reading the step diagnostics of the quartic saddle notebooks.
%
% Companion to notebooks/Saddle/step_diagnostics_v1.ipynb. Every number in this
% file is printed by a cell of that notebook that ran; none is transcribed from
% anywhere else.
%
% Standalone. Own preamble, no bibliography, no \cite. The macro block mirrors
% the one near the top of _Thesis_FINAL/Article3/Survey-part3-v1.tex, through
% \providecommand, so this file can later be folded into that document.
% =============================================================================
\documentclass[11pt,a4paper]{article}

\usepackage[utf8]{inputenc}
\usepackage[british]{babel}
\usepackage{amsmath, amssymb, amsthm}
\usepackage{booktabs}
\usepackage{graphicx}
\usepackage{geometry}
\usepackage{float}
\usepackage{enumitem}
\usepackage[hidelinks]{hyperref}
\geometry{margin=2.4cm}

% ---- macros mirrored from Survey-part3-v1.tex -------------------------------
\providecommand{\defeq}{\mathrel{:=}}
\providecommand{\kbar}{\overline{\kappa}}
\providecommand{\calO}{{\mathcal{O}}}
\providecommand{\calK}{{\mathcal{K}}}
\providecommand{\NN}{{\mathbb{N}}}
\providecommand{\abs}[1]{\left \vert #1 \right \vert}
\providecommand{\norm}[1]{\left \lVert #1 \right \rVert}
\providecommand*\mykappa[1]{ \kappa_{\scriptscriptstyle \text{#1}} }
\providecommand*\myeta[1]{ \eta_{\scriptscriptstyle \text{#1}} }
\providecommand{\Rdelta}{\textsf{R-delta}}
\providecommand{\Rstep}{\textsf{R-step}}
\providecommand{\Rdfo}{\textsf{R-DFO}}
\providecommand{\Rgrad}{\textsf{R-grad}}
\providecommand{\Rrtr}{\textsf{R-RTR}}
\providecommand{\Rgrtr}{\textsf{R-gRTR}}
\providecommand{\gone}{\gamma_1}
\providecommand{\gtwo}{\gamma_2}
\providecommand{\gthree}{\gamma_3}

\newcommand{\e}[1]{\times 10^{#1}}
\newtheorem{proposition}{Proposition}
\theoremstyle{remark}
\newtheorem{remark}{Remark}

\title{Reading the step diagnostics\\
       \large A companion to \texttt{step\_diagnostics\_v1.ipynb}}
\author{}
\date{}

\begin{document}
\maketitle

\begin{abstract}
We explain the eight columns of the step-diagnostic tables and the six functions
that produce them, and we interpret the tables of
\texttt{step\_diagnostics\_v1.ipynb}. Truncated conjugate gradient leaves its loop
in three ways and every one of them returns a positive multiple of $-g_k$, so a
one-dimensional Krylov subspace forces the cosine test to pass. The two Cauchy
statistics are therefore nested rather than independent, $\texttt{cg1} \leq
\texttt{cauchy}$, which is why they print the same number on $28$ of $30$ rows.
On the two rows where they separate the accepted explanation, that $g_k$ is near
an eigenvector, is refuted by the measurement. The quantity that separates the
populations is the gap between the first conjugate-gradient iterate and the
boundary, whose median is $6.5\times10^{-9}$ at the disagreeing iterations and
$1.6\times10^{-3}$ at the agreeing ones. We give the arithmetic that turns that
gap into an angle below the $10^{-12}$ threshold, and we check it. The truncation
condition is confirmed on $17755$ of $17763$ exact-Hessian iterations, and all
eight exceptions are the one conjugate-gradient exit the condition does not
model.
\end{abstract}

% =============================================================================
\section{The setting}\label{sec: setting}
% =============================================================================

The test function is
\begin{equation}
  f(p) = p_1^4 - p_1^3
       + \left(\tfrac14 - \tfrac{p_1}{2}\right)p_2^2
       + \tfrac{p_2^4}{4},
  \label{eqn: f}
\end{equation}
with four critical points. The origin is degenerate, $\nabla^2 f(0,0)$ being
singular and positive semidefinite. The point $(0.75, 0)$ is a strict saddle with
$\lambda_{\min} = -1/4$. Two local minimisers sit at
$\bigl(\tfrac{1+\sqrt5}{4}, \pm\sqrt{\tfrac{\sqrt5-1}{4}}\bigr)$.

Every run in the notebook uses $\Rgrad$ capped, so
$\Delta_k = \mu_k\norm{g_k}$ with $\mu_k \leq \overline{\mu}$, sweeping
$\overline{\mu}$ over
$\{0.01, 0.05, 0.1, 0.3, 0.5, 1, 2, 4, 8, 128\}$ against three model Hessians,
\texttt{ExactHessian}, \texttt{SR1} and \texttt{LBFGS}. The subproblem solver is
truncated conjugate gradient throughout, apart from one comparison against
\texttt{ExactMS}. The starting point is $(-0.5, 0.6)$ and the tolerance
$10^{-9}$, with at most $2000$ iterations.

Everything below is read from \texttt{stats.solver\_specific}. Nothing is
recomputed from the iterates that the trace already carries.

\subsection{The three ways truncated conjugate gradient stops}
\label{sec: exits}

The whole of Section~\ref{sec: columns} follows from the recurrence, so it is
worth setting down first. \texttt{SteihaugCG} starts at $s = 0$ with residual
$r_0 = -g_k$ and first direction $d_0 = -g_k$, and it leaves the loop in exactly
three ways.

\begin{enumerate}[label=(E\arabic*), leftmargin=3em]
  \item \label{it: negcurv}
    \emph{Negative curvature.} If $d_0^{\!\top}B_kd_0 \leq 0$ the recurrence runs
    to the boundary along $d_0$ and returns $s_k = -\Delta_k g_k/\norm{g_k}$.
  \item \label{it: boundary}
    \emph{The boundary.} Otherwise it forms
    $\alpha_0 = \norm{g_k}^2/(g_k^{\!\top}B_kg_k)$ and the candidate
    $z_1 = -\alpha_0 g_k$. If $\norm{z_1} \geq \Delta_k$ it truncates to the
    boundary and returns $s_k = -\Delta_k g_k/\norm{g_k}$.
  \item \label{it: forcing}
    \emph{The forcing condition.} Otherwise it accepts $z_1$, and if the new
    residual satisfies
    $\norm{r_1} \leq \min\{\chi, \norm{g_k}^\theta\}\norm{g_k}$ it stops there,
    returning $s_k = z_1 = -\alpha_0 g_k$.
\end{enumerate}

Only if none of the three fires does the recurrence build a second direction.

The consequence used throughout is immediate. \emph{In all three exits the
returned step is a positive multiple of $-g_k$.} A run that leaves after one
conjugate-gradient iteration has taken a step along the steepest descent
direction, whatever the model Hessian was.

% =============================================================================
\section{The six columns of Table 1}\label{sec: columns}
% =============================================================================

\subsection{\texttt{trunc}, \texttt{inside} and \texttt{krylov}}

These are the three step classes of \texttt{step\_classes}, as fractions of the
iterations that happened. Write $c_k$ for the inner iteration count and $a_k$ for
the activity flag, which is true when $\norm{s_k} = \Delta_k$.
\begin{align}
  \texttt{trunc}  &= \tfrac1K \,\#\{k : c_k = 1 \text{ and } a_k\},
  \label{eqn: trunc} \\
  \texttt{inside} &= \tfrac1K \,\#\{k : c_k = 1 \text{ and } \lnot a_k\},
  \label{eqn: inside} \\
  \texttt{krylov} &= \tfrac1K \,\#\{k : c_k \geq 2\}.
  \label{eqn: krylov}
\end{align}
The three partition the iterations, so they sum to one.

\texttt{trunc} collects \ref{it: negcurv} and \ref{it: boundary}. The step is
$-\Delta_k g_k/\norm{g_k}$, so the model Hessian has touched neither the
direction nor the length: the direction is $-g_k$ and the length is $\Delta_k$.
\textbf{This is the only one of the three that is gradient descent in disguise},
and the step length is the radius rather than anything the model chose.

\texttt{inside} is \ref{it: forcing}. The direction is again $-g_k$, but the
length is $\alpha_0\norm{g_k} = \norm{g_k}^3/(g_k^{\!\top}B_kg_k)$, which is the
exact minimiser of the model along $-g_k$. The Hessian set the length and not the
direction. This is the Cauchy point in the strict sense.

\texttt{krylov} is everything else, where a second direction entered and the
Hessian influenced the direction as well.

\emph{Interpretation.} Read \texttt{trunc} as the fraction of the run on which
the trust region, rather than the model, decided the step. A \texttt{trunc} of
one means the method was gradient descent with a prescribed step length for the
whole run. \texttt{krylov} near one means the model was doing the work.

\texttt{step\_classes} returns $(\mathrm{NaN}, \mathrm{NaN}, \mathrm{NaN})$ under
\texttt{ExactMS}, which reports $c_k = 0$ throughout, since a direct solve has no
inner iteration count. The notebook confirms this: \texttt{cg\_iters} takes the
values $\{2, 1\}$ under \texttt{SteihaugCG} and $\{0\}$ under \texttt{ExactMS},
and the classes are $(0.999, 0.000, 0.001)$ and
$(\mathrm{NaN}, \mathrm{NaN}, \mathrm{NaN})$ respectively.

\subsection{\texttt{cauchy}}

The cosine-based fraction reported by \texttt{hypotheses\_report},
\begin{equation}
  \texttt{cauchy}
    = \tfrac1K \,\#\Bigl\{k : \bigl\lvert \cos(s_k, -g_k) - 1 \bigr\rvert
                                \leq 10^{-12}\Bigr\},
  \qquad
  \cos(s_k,-g_k) = \frac{-g_k^{\!\top}s_k}{\norm{g_k}\,\norm{s_k}}.
  \label{eqn: cauchy}
\end{equation}
It measures how often the step \emph{pointed} along steepest descent, to within
$10^{-12}$. It says nothing about the length.

\subsection{\texttt{cg1}}

The Krylov dimension,
\begin{equation}
  \texttt{cg1} = \tfrac1K \,\#\{k : c_k = 1\}
              = \texttt{trunc} + \texttt{inside}.
  \label{eqn: cg1}
\end{equation}
It measures how often the recurrence stopped before building a second direction.
The identity with the two step classes is exact, not approximate, since
\eqref{eqn: trunc} and \eqref{eqn: inside} split $\{c_k = 1\}$ on the activity
flag.

\subsection{\texttt{cos\_min}}

The minimum of $\cos(s_k,-g_k)$ over the run,
\begin{equation}
  \texttt{cos\_min} = \min_{k} \cos(s_k, -g_k).
  \label{eqn: cosmin}
\end{equation}
This is the column that stops a fraction from being read as a description of the
whole run. \texttt{cauchy} says how \emph{often} the step was Cauchy;
\texttt{cos\_min} says how far from Cauchy the run ever got. A run with
$\texttt{cos\_min} = 1$ was gradient descent at \emph{every} iteration, not at
most of them.

\emph{Read \texttt{cos\_min} rather than the cosine at the final iterate.} The
last value is one number out of however many iterations ran, and the two
disagree: in the companion notebook \texttt{exact\_steihaug\_v1} the final-iterate
cosine is $1.00000$ at $\overline{\mu} = 4$ while \texttt{cauchy} is $0.4885$
there.

% =============================================================================
\section{Why \texttt{cauchy} and \texttt{cg1} answer the same question}
\label{sec: same question}
% =============================================================================

Both ask whether the model Hessian influenced the \emph{direction} of the step.
They ask it of different objects. \texttt{cg1} asks it of the Krylov subspace,
$\dim\calK_k = 1$, and \texttt{cauchy} asks it of the angle the step came out at.

The two are linked in one direction by Section~\ref{sec: exits}.

\begin{proposition}[the inclusion]\label{prop: inclusion}
  For every $k$ with $c_k = 1$, the returned step satisfies $s_k = -t g_k$ for
  some $t > 0$, hence $\cos(s_k,-g_k) = 1$ exactly. Consequently
  \begin{equation}
    \{k : c_k = 1\} \;\subseteq\;
    \bigl\{k : \abs{\cos(s_k,-g_k) - 1} \leq 10^{-12}\bigr\},
    \qquad\text{so}\qquad
    \texttt{cg1} \;\leq\; \texttt{cauchy}.
    \label{eqn: inclusion}
  \end{equation}
\end{proposition}

\begin{proof}
  Each of \ref{it: negcurv}, \ref{it: boundary} and \ref{it: forcing} returns a
  positive multiple of $d_0 = -g_k$, the first with $t = \Delta_k/\norm{g_k}$,
  the second likewise, the third with $t = \alpha_0 > 0$.
\end{proof}

So \texttt{cauchy} counts everything \texttt{cg1} counts, and possibly more. The
inequality \eqref{eqn: inclusion} is the whole relationship, and the two agree
whenever the extra set is empty.

\subsection{Why they hold the same value on nearly every row}

Equality fails only at an iteration with $c_k \geq 2$ whose returned step is
nonetheless within $10^{-12}$ of the $-g_k$ direction. That asks the second and
later conjugate-gradient directions to add nothing \emph{measurable} to the
direction, while still being taken.

That is a demanding requirement, and on this problem it is met on $2$ of the $30$
rows of Table~\ref{tab: one}. On the remaining $28$ the extra set is empty and the
two columns print the same number, which is not a coincidence but
Proposition~\ref{prop: inclusion} with nothing left over.

\subsection{The two rows where they differ, and why}

Under \texttt{ExactHessian} at $\overline{\mu} = 4$ and $\overline{\mu} = 8$ the
columns separate, $0.4885$ against $0.3345$ and $0.4815$ against $0.3319$. Every
one of the $308$ and $251$ disagreeing iterations goes the same way, the cosine
counting an iteration the Krylov dimension does not, exactly as
\eqref{eqn: inclusion} requires.

The natural explanation is that $g_k$ is close to an eigenvector of $H_k$ there,
so the second direction adds nothing. \textbf{It does not survive the
measurement.} Pooled over the disagreeing runs, \texttt{eig\_deviation} has median
$7.0695 \times 10^{-1}$ at the disagreeing iterations against
$7.0617 \times 10^{-1}$ at the agreeing ones. The two populations are
indistinguishable on that statistic, and both are far from zero.

What separates them is the boundary. Writing
$\varepsilon_k = 1 - \norm{z_1}/\Delta_k$ for the gap between the first
conjugate-gradient iterate and the trust-region boundary, the median of
$\varepsilon_k$ is $6.536\times10^{-9}$ at the disagreeing iterations against
$1.6266\times10^{-3}$ at the agreeing ones, five orders of magnitude apart.

The mechanism follows. At those iterations $\norm{z_1}$ falls a relative
$6.5\times10^{-9}$ short of $\Delta_k$, so \ref{it: boundary} does not fire and
the recurrence continues. The radial room left is only about
$\varepsilon_k\Delta_k$. Moving from $z_1$ to the boundary along any direction
making an angle $\varphi$ with $z_1$ costs a distance $\varepsilon_k\Delta_k /
\cos\varphi$, so the step can rotate by at most about $\varepsilon_k$ radians,
and
\begin{equation}
  1 - \cos(s_k,-g_k) \;\lesssim\; \tfrac12\varepsilon_k^2 .
  \label{eqn: gap arithmetic}
\end{equation}

The notebook measures both sides. With the median $\varepsilon_k =
6.536\times10^{-9}$ the bound \eqref{eqn: gap arithmetic} gives
$2.136\times10^{-17}$, and the measured median of $\abs{1 - \cos(s_k,-g_k)}$ at
the disagreeing iterations is $2.220\times10^{-16}$, which is
$\texttt{eps(Float64)}$ exactly. The geometric bound sits below the rounding
floor, so the measurement is floor limited rather than geometry limited. The
measured value is four orders of magnitude below the $10^{-12}$ threshold of
\eqref{eqn: cauchy} and the bound is five. The Krylov dimension counts two and
the angle cannot tell.

The threshold is nonetheless doing real work at the tail. The \emph{largest}
$\abs{1 - \cos(s_k,-g_k)}$ among the disagreeing iterations is
$9.2293\times10^{-13}$, that is $92\%$ of $10^{-12}$. The corresponding figure
over the agreeing iterations has median $1.3101\times10^{-6}$ and maximum
$3.2434\times10^{-1}$. The two populations are separated by ten orders of
magnitude in the median, and one disagreeing iteration nonetheless sits close to
the boundary of the test.

\begin{remark}[which one to use]\label{rem: which}
  \textbf{Use \texttt{cg1} as the definition of truncation and the cosine as
  corroboration.} The truncation condition is a statement about the
  conjugate-gradient recurrence, so the Krylov dimension is what states it, and
  the angle is a consequence of it by Proposition~\ref{prop: inclusion}. Where
  they disagree, \eqref{eqn: gap arithmetic} shows the angle is reporting an
  arithmetic near-coincidence rather than a property of the step.
\end{remark}

% =============================================================================
\section{The six diagnostics}\label{sec: six}
% =============================================================================

All six live in \texttt{saddle\_problem.jl} and are shared by the four saddle
notebooks.

\paragraph{\texttt{step\_classes(ss)}.}
The triple \eqref{eqn: trunc} to \eqref{eqn: krylov}, read from
\texttt{:cg\_iters\_trajectory} and \texttt{:active\_trajectory}.
\emph{Does not apply} under \texttt{ExactMS}, which reports $c_k = 0$, and
returns $(\mathrm{NaN}, \mathrm{NaN}, \mathrm{NaN})$ there rather than a fraction
of zero. Read the angle instead in that case.

\paragraph{\texttt{eig\_deviation(x)}.}
The sine of the angle between $H(x)g(x)$ and $g(x)$,
\begin{equation}
  \texttt{eig\_deviation}(x)
    = \sqrt{1 - \left(\frac{\abs{(Hg)^{\!\top}g}}{\norm{Hg}\,\norm{g}}\right)^{\!2}},
  \label{eqn: eigdev}
\end{equation}
zero exactly when $g$ is an eigenvector of $H$. On \eqref{eqn: f} the Hessian is
diagonal on the axis $p_2 = 0$ and the gradient lies along $e_1$ there, so the
deviation vanishes on the axis. The notebook confirms it is $0$ to the last bit
at $(0.3, 0)$, $(0.6, 0)$ and $(-0.5, 0)$, and $2.415\times10^{-1}$ at the
starting point. \emph{Does not apply} where $g = 0$ or $Hg = 0$, returning
\texttt{NaN}.

\paragraph{\texttt{z1\_ratio(x, $\Delta$)}.}
The first Steihaug iterate against the radius,
\begin{equation}
  \texttt{z1\_ratio}(x, \Delta)
    = \frac{\norm{z_1}}{\Delta}
    = \frac{\norm{g}^{3}}{\bigl(g^{\!\top}Hg\bigr)\,\Delta}.
  \label{eqn: z1}
\end{equation}
It reaches one exactly when \ref{it: boundary} fires, so the truncation condition
is
\begin{equation}
  \Delta \;\leq\; \frac{\norm{g}^{3}}{g^{\!\top}Hg},
  \qquad\text{equivalently}\qquad
  \mu \;\leq\; \frac{\norm{g}^{2}}{g^{\!\top}Hg}
  \ \text{ when } \Delta = \mu\norm{g},
  \label{eqn: trunc condition}
\end{equation}
a reciprocal curvature. A cap below that quantity truncates every step, which is
the mechanism behind the top rows of Table~\ref{tab: one}.
\emph{Does not apply} where $g^{\!\top}Hg \leq 0$, returning \texttt{Inf},
because conjugate gradient then goes straight to the boundary through
\ref{it: negcurv}. The notebook exhibits six such points, all near the axis
between the origin and the saddle, with $g^{\!\top}Hg$ from
$-4.33\times10^{-8}$ to $-2.15\times10^{-3}$.

\paragraph{\texttt{spectrum\_variance(x)}.}
The $g$-weighted variance of the spectrum of $H(x)$. With $H = V\Lambda V^{\!\top}$
and weights $w_i = (V^{\!\top}g)_i^2/\norm{g}^2$,
\begin{equation}
  \texttt{spectrum\_variance}(x)
    = \sum_i w_i\lambda_i^2 - \Bigl(\sum_i w_i\lambda_i\Bigr)^{\!2}.
  \label{eqn: specvar}
\end{equation}
It is the coefficient in the \texttt{ExactMS} angle law of the first notebook,
$1 - \cos(s_k,-g_k) \approx \tfrac12\,\texttt{spectrum\_variance}\,(\Delta_k/\norm{g_k})^2$.
It vanishes where $g$ is an eigenvector, so it is $0$ on the axis $p_2 = 0$,
which the notebook shows beside \eqref{eqn: eigdev}. \emph{Does not apply} where
$g = 0$, returning \texttt{NaN}.

\paragraph{\texttt{loglog\_fit(x, y)}.}
Least squares slope of $\log_{10}y$ on $\log_{10}x$, with intercept, $R^2$ and
the number of usable points. Points with $x \leq 0$, $y \leq 0$ or a non-finite
value are dropped first. \emph{Does not apply} below five usable points, where it
returns \texttt{NaN}, so a slope is never quoted from a handful of iterations.
The notebook shows \texttt{NaN} at three and four points and slope $1.5$ with
$R^2 = 1.0$ at five and eight. \textbf{Report $n$ and $R^2$ beside every slope
taken from it.}

\paragraph{\texttt{nearest\_crit(x)}.}
The nearest of the four critical points and its distance, as a pair.
\texttt{which\_crit} is written in terms of it, returning the name when the
distance is below its tolerance and \texttt{"other"} otherwise, so a limit point
is classified by one distance computation. The notebook shows the two agreeing at
the four critical points and separating at $(-0.5, 0.6)$, where the nearest is
the origin at distance $7.810\times10^{-1}$ and \texttt{which\_crit} returns
\texttt{"other"}.

\emph{Never report \texttt{"other"} as an outcome without the distance and
$\norm{g}$ beside it.} A run that exhausted its budget while still moving is not
a run that converged somewhere else, and only those two numbers separate the
cases.

% =============================================================================
\section{The tables}\label{sec: tables}
% =============================================================================

\subsection{Table 1: the degeneration ladder}\label{sec: table one}

\begin{table}[H]
  \centering\small
  \caption{$\Rgrad$ capped, truncated conjugate gradient, three model Hessians.
    The columns are \eqref{eqn: trunc} to \eqref{eqn: cosmin}. Note
    $\texttt{cg1} = \texttt{trunc} + \texttt{inside}$ on every row, to the
    printed precision, which is \eqref{eqn: cg1}.}
  \label{tab: one}
  \begin{tabular}{llrrrrrrrr}
    \toprule
    model & $\overline{\mu}$ & status & iter
      & \texttt{trunc} & \texttt{inside} & \texttt{krylov}
      & \texttt{cauchy} & \texttt{cg1} & \texttt{cos\_min} \\
    \midrule
    ExactHessian & 0.01 & max\_iter & 2000 & 1.000 & 0.000 & 0.000 & 1.0000 & 1.0000 & 1.000000 \\
    ExactHessian & 0.05 & max\_iter & 2000 & 1.000 & 0.000 & 0.000 & 1.0000 & 1.0000 & 1.000000 \\
    ExactHessian & 0.10 & max\_iter & 2000 & 1.000 & 0.000 & 0.000 & 1.0000 & 1.0000 & 1.000000 \\
    ExactHessian & 0.30 & max\_iter & 2000 & 1.000 & 0.000 & 0.001 & 0.9995 & 0.9995 & 0.877285 \\
    ExactHessian & 0.50 & max\_iter & 2000 & 0.999 & 0.000 & 0.001 & 0.9990 & 0.9990 & 0.877285 \\
    ExactHessian & 1.00 & max\_iter & 2000 & 0.999 & 0.000 & 0.002 & 0.9985 & 0.9985 & 0.877285 \\
    ExactHessian & 2.00 & max\_iter & 2000 & 0.998 & 0.001 & 0.002 & 0.9985 & 0.9985 & 0.877285 \\
    ExactHessian & 4.00 & max\_iter & 2000 & 0.334 & 0.001 & 0.665 & 0.4885 & 0.3345 & 0.877285 \\
    ExactHessian & 8.00 & first\_order & 1678 & 0.331 & 0.001 & 0.668 & 0.4815 & 0.3319 & 0.675664 \\
    ExactHessian & 128.00 & first\_order & 85 & 0.047 & 0.035 & 0.918 & 0.0824 & 0.0824 & 0.246457 \\
    \midrule
    SR1 & 0.01 & max\_iter & 2000 & 1.000 & 0.000 & 0.000 & 1.0000 & 1.0000 & 1.000000 \\
    SR1 & 0.05 & max\_iter & 2000 & 1.000 & 0.000 & 0.000 & 1.0000 & 1.0000 & 1.000000 \\
    SR1 & 0.10 & max\_iter & 2000 & 1.000 & 0.000 & 0.000 & 1.0000 & 1.0000 & 1.000000 \\
    SR1 & 0.30 & max\_iter & 2000 & 1.000 & 0.000 & 0.000 & 1.0000 & 1.0000 & 1.000000 \\
    SR1 & 0.50 & stalled & 83 & 1.000 & 0.000 & 0.000 & 1.0000 & 1.0000 & 1.000000 \\
    SR1 & 1.00 & first\_order & 23 & 0.739 & 0.000 & 0.261 & 0.7391 & 0.7391 & 0.671198 \\
    SR1 & 2.00 & first\_order & 24 & 0.625 & 0.042 & 0.333 & 0.6667 & 0.6667 & 0.692186 \\
    SR1 & 4.00 & first\_order & 19 & 0.632 & 0.105 & 0.263 & 0.7368 & 0.7368 & 0.761860 \\
    SR1 & 8.00 & first\_order & 19 & 0.632 & 0.105 & 0.263 & 0.7368 & 0.7368 & 0.761860 \\
    SR1 & 128.00 & first\_order & 19 & 0.632 & 0.105 & 0.263 & 0.7368 & 0.7368 & 0.761860 \\
    \midrule
    LBFGS & 0.01 & max\_iter & 2000 & 1.000 & 0.000 & 0.000 & 1.0000 & 1.0000 & 1.000000 \\
    LBFGS & 0.05 & max\_iter & 2000 & 1.000 & 0.000 & 0.000 & 1.0000 & 1.0000 & 1.000000 \\
    LBFGS & 0.10 & max\_iter & 2000 & 1.000 & 0.000 & 0.000 & 1.0000 & 1.0000 & 1.000000 \\
    LBFGS & 0.30 & max\_iter & 2000 & 1.000 & 0.000 & 0.000 & 1.0000 & 1.0000 & 1.000000 \\
    LBFGS & 0.50 & stalled & 85 & 1.000 & 0.000 & 0.000 & 1.0000 & 1.0000 & 1.000000 \\
    LBFGS & 1.00 & first\_order & 30 & 0.333 & 0.100 & 0.567 & 0.4333 & 0.4333 & 0.661424 \\
    LBFGS & 2.00 & stalled & 28 & 0.464 & 0.143 & 0.393 & 0.6071 & 0.6071 & 0.599361 \\
    LBFGS & 4.00 & first\_order & 21 & 0.333 & 0.238 & 0.429 & 0.5714 & 0.5714 & 0.426056 \\
    LBFGS & 8.00 & first\_order & 21 & 0.333 & 0.238 & 0.429 & 0.5714 & 0.5714 & 0.426056 \\
    LBFGS & 128.00 & first\_order & 21 & 0.333 & 0.238 & 0.429 & 0.5714 & 0.5714 & 0.426056 \\
    \bottomrule
  \end{tabular}
\end{table}

Read the table down each block. Three regimes appear, and the cap is what moves
between them.

\paragraph{The degenerate regime, $\overline{\mu} \leq 0.1$.}
$\texttt{trunc} = 1.000$, $\texttt{cauchy} = \texttt{cg1} = 1.0000$ and
$\texttt{cos\_min} = 1.000000$, on all three models. Every step of two thousand
was truncated at the boundary along $-g_k$. The method was gradient descent with
step length $\Delta_k$ for the whole run, and $\texttt{cos\_min}$ is what says
\emph{whole}: a fraction of one leaves open a single stray iteration, and a
minimum of one does not.

The mechanism is \eqref{eqn: trunc condition}. A cap below
$\norm{g_k}^2/(g_k^{\!\top}H_kg_k)$ truncates every step, and the model never
enters. That the three models agree exactly here is the point. When the region is
this small the model Hessian is irrelevant, so \texttt{SR1} and \texttt{LBFGS}
produce the same run as the exact Hessian.

\paragraph{The mixed regime.}
Between $\overline{\mu} = 0.3$ and $2$ under \texttt{ExactHessian}, \texttt{trunc}
stays above $0.998$ while \texttt{cos\_min} drops to $0.877285$. The fractions
barely move and the minimum does, which is the case the two kinds of column exist
to separate. One or two iterations in two thousand escaped the truncated branch,
and only \texttt{cos\_min} records that they did.

\paragraph{The Krylov regime.}
At $\overline{\mu} = 128$ under \texttt{ExactHessian}, $\texttt{krylov} = 0.918$
and the run converges in $85$ iterations rather than exhausting $2000$.
$\texttt{cos\_min} = 0.246457$, so the step direction departed from $-g_k$ by
about $76^\circ$ at its worst. The model is choosing the direction, and the run
is a trust-region method rather than a disguised gradient method.

\paragraph{The quasi-Newton models reach that regime sooner.}
\texttt{SR1} and \texttt{LBFGS} leave the degenerate regime at
$\overline{\mu} = 1$, where \texttt{ExactHessian} is still at
$\texttt{trunc} = 0.999$. Their runs then take $19$ to $30$ iterations against
$2000$. The cap that degenerates a method is not a property of the cap alone.

\paragraph{\texttt{inside} is almost never occupied.}
It is $0.000$ or $0.001$ on every \texttt{ExactHessian} row and reaches $0.238$
only under \texttt{LBFGS}. The degeneration proceeds by truncation rather than by
a short interior step, which is why \texttt{trunc} and not \texttt{cg1} is the
column that identifies gradient descent.

\subsection{Table 2: the two Cauchy statistics}\label{sec: table two}

Table~\ref{tab: one} already shows the answer to the question of
Section~\ref{sec: same question}. On $28$ of the $30$ rows \texttt{cauchy} and
\texttt{cg1} print the same number, which is Proposition~\ref{prop: inclusion}
with an empty remainder. On the two \texttt{ExactHessian} rows at
$\overline{\mu} = 4$ and $8$ they separate, and the notebook's Table 2 counts
$308$ and $251$ disagreeing iterations, all in the direction the inclusion
allows.

\begin{table}[H]
  \centering\small
  \caption{The two candidate explanations, pooled over the runs that disagree at
    all. \texttt{eig\_deviation} does not separate the two populations and the
    boundary gap $1 - \norm{z_1}/\Delta$ separates them by five orders of
    magnitude.}
  \label{tab: two}
  \begin{tabular}{lrrr}
    \toprule
    quantity & min & median & max \\
    \midrule
    \texttt{eig\_deviation}, disagreeing & $7.0343\e{-1}$ & $7.0695\e{-1}$ & $8.6594\e{-1}$ \\
    \texttt{eig\_deviation}, agreeing    & $7.5345\e{-3}$ & $7.0617\e{-1}$ & $8.6594\e{-1}$ \\
    \midrule
    $1 - \norm{z_1}/\Delta$, disagreeing & $4.4409\e{-16}$ & $6.5360\e{-9}$ & $2.2238\e{-6}$ \\
    $1 - \norm{z_1}/\Delta$, agreeing    & $-1.0790\e{0}$  & $1.6266\e{-3}$ & $6.8788\e{-1}$ \\
    \bottomrule
  \end{tabular}
\end{table}

The negative entries in the last row are not an error. On a truncating iteration
$\norm{z_1} \geq \Delta$, so $1 - \norm{z_1}/\Delta \leq 0$, and the agreeing
population is dominated by such iterations. The disagreeing population sits at
$+6.5\times10^{-9}$, that is just \emph{inside} the boundary, which is exactly
the configuration \ref{it: boundary} fails to catch.

\subsection{Table 3: the truncation condition, iteration by iteration}
\label{sec: table three}

\begin{table}[H]
  \centering\small
  \caption{The prediction $\norm{z_1}/\Delta_k \geq 1$ against the observed
    $c_k = 1$. Rows with no disagreement under \texttt{SR1} and \texttt{LBFGS}
    are omitted; the counts below are over all ten caps.}
  \label{tab: three}
  \begin{tabular}{llrrrr}
    \toprule
    model & $\overline{\mu}$ & iterations & predicted & observed & disagreements \\
    \midrule
    ExactHessian & 0.01 & 2000 & 2000 & 2000 & 0 \\
    ExactHessian & 0.05 & 2000 & 2000 & 2000 & 0 \\
    ExactHessian & 0.10 & 2000 & 2000 & 2000 & 0 \\
    ExactHessian & 0.30 & 2000 & 1999 & 1999 & 0 \\
    ExactHessian & 0.50 & 2000 & 1998 & 1998 & 0 \\
    ExactHessian & 1.00 & 2000 & 1997 & 1997 & 0 \\
    ExactHessian & 2.00 & 2000 & 1996 & 1997 & 1 \\
    ExactHessian & 4.00 & 2000 & 667 & 669 & 2 \\
    ExactHessian & 8.00 & 1678 & 555 & 557 & 2 \\
    ExactHessian & 128.00 & 85 & 4 & 7 & 3 \\
    \midrule
    SR1 & 0.30 & 2000 & 1999 & 2000 & 1 \\
    SR1 & 0.50 & 82 & 81 & 82 & 1 \\
    SR1 & 1.00 & 23 & 16 & 17 & 1 \\
    SR1 & 2.00 & 24 & 14 & 16 & 2 \\
    SR1 & 4.00 & 19 & 11 & 14 & 3 \\
    \midrule
    LBFGS & 0.30 & 2000 & 1999 & 2000 & 1 \\
    LBFGS & 0.50 & 84 & 83 & 84 & 1 \\
    LBFGS & 1.00 & 30 & 19 & 13 & 12 \\
    LBFGS & 2.00 & 27 & 16 & 16 & 4 \\
    LBFGS & 4.00 & 21 & 11 & 12 & 5 \\
    \bottomrule
  \end{tabular}
\end{table}

Under \texttt{ExactHessian} the prediction is right on $17755$ of $17763$
iterations, a disagreement rate of $0.0005$. The rate is $0.0017$ under
\texttt{SR1} and $0.0040$ under \texttt{LBFGS}, and that difference is a property
of the diagnostic rather than of the condition: \texttt{z1\_ratio} reads the
\emph{true} Hessian through \texttt{hess}, and conjugate gradient runs on the
\emph{model} Hessian. The two coincide under \texttt{ExactHessian} and nowhere
else.

\begin{table}[H]
  \centering\small
  \caption{Every disagreement under \texttt{ExactHessian}, individually. Three
    distinct iterates account for all eight, appearing once per cap large enough
    to reach them.}
  \label{tab: four}
  \begin{tabular}{lrrrrrr}
    \toprule
    $\overline{\mu}$ & $k$ & \texttt{z1\_ratio} & $c_k$
      & \texttt{eig\_deviation} & $\Delta_k$ & $\norm{g_k}$ \\
    \midrule
    2.00 & 3 & 0.74290 & 1 & 7.5345e-03 & 6.1095e-02 & 3.0547e-02 \\
    4.00 & 3 & 0.74290 & 1 & 7.5345e-03 & 6.1095e-02 & 3.0547e-02 \\
    4.00 & 4 & 0.76071 & 1 & 5.3537e-02 & 3.1586e-02 & 7.8966e-03 \\
    8.00 & 3 & 0.74290 & 1 & 7.5345e-03 & 6.1095e-02 & 3.0547e-02 \\
    8.00 & 4 & 0.76071 & 1 & 5.3537e-02 & 3.1586e-02 & 7.8966e-03 \\
    128.00 & 3 & 0.74290 & 1 & 7.5345e-03 & 6.1095e-02 & 3.0547e-02 \\
    128.00 & 4 & 0.76071 & 1 & 5.3537e-02 & 3.1586e-02 & 7.8966e-03 \\
    128.00 & 7 & 0.78195 & 1 & 9.8707e-03 & 4.1855e-03 & 1.3080e-04 \\
    \bottomrule
  \end{tabular}
\end{table}

All eight have the same shape. \texttt{z1\_ratio} lies between $0.7429$ and
$0.7819$, so the prediction says the step will not truncate, and $c_k = 1$ says
it stopped after one iteration anyway. None is within $10^{-6}$ of one and none
is \texttt{Inf}, so neither \ref{it: boundary} nor \ref{it: negcurv} explains
them.

The cause is \ref{it: forcing}. \texttt{z1\_ratio} models the boundary exit and
nothing else, while the recurrence has three exits. Those eight iterations
stopped \emph{inside} the region because the residual met the forcing condition,
which \eqref{eqn: z1} cannot see.

\begin{remark}[the prediction is one-sided, and that is what to rely on]
\label{rem: onesided}
  There is no disagreement in the other direction. Over all $17763$
  \texttt{ExactHessian} iterations, \texttt{z1\_ratio} $\geq 1$ never occurs with
  $c_k \geq 2$. So $\norm{z_1}/\Delta_k \geq 1$ is \emph{sufficient} for
  truncation and not necessary, which is what \eqref{eqn: trunc condition}
  asserts and all that Section~\ref{sec: table one} needs from it.
\end{remark}

\subsection{What the four tables say together}\label{sec: together}

The cap controls the radius, the radius controls the truncation through
\eqref{eqn: trunc condition}, and the truncation controls whether the model
Hessian influences the step at all. Table~\ref{tab: one} measures the last of
those three, Table~\ref{tab: three} confirms the middle one to a rate of five in
ten thousand, and Tables~\ref{tab: two} and \ref{tab: four} account for every
iteration where the diagnostics disagree.

Two of the eight columns are load bearing and the rest corroborate.
\texttt{trunc} identifies gradient descent in disguise, because it is the only
class in which neither the direction nor the length came from the model.
\texttt{cos\_min} says whether that held throughout the run or only on average.
\texttt{cauchy}, \texttt{cg1}, \texttt{inside} and \texttt{krylov} are consistent
with those two by \eqref{eqn: cg1} and Proposition~\ref{prop: inclusion}, and
where they part company the reason is arithmetic rather than algorithmic.

\end{document}
